
\documentclass{article}
\usepackage{colortbl}
\usepackage{makecell}
\usepackage{multirow}
\usepackage{supertabular}

\begin{document}

\newcounter{utterance}

\centering \large Interaction Transcript for game `eqbench', experiment `eq\_bench', episode 111 with reflexion{-}llama{-}3.1{-}8b{-}instruct.
\vspace{24pt}

{ \footnotesize  \setcounter{utterance}{1}
\setlength{\tabcolsep}{0pt}
\begin{supertabular}{c@{$\;$}|p{.15\linewidth}@{}p{.15\linewidth}p{.15\linewidth}p{.15\linewidth}p{.15\linewidth}p{.15\linewidth}}
    \# & \multicolumn{2}{c}{Player} && \multicolumn{2}{c}{Game Master} \\
    \hline

    \theutterance \stepcounter{utterance}  
    & & & \multicolumn{4}{p{0.6\linewidth}}{
        \cellcolor[rgb]{0.9,0.9,0.9}{
            \makecell[{{p{\linewidth}}}]{
                \texttt{\tiny{[P1$\langle$GM]}}
                \texttt{You have attempted this same task before and did not succeed. These are the plans you wrote after those attempts.} \\
\texttt{Trial \#1: New plan: } \\
\\ 
\texttt{Upon re{-}reading the dialogue, I realized that I focused too much on the surface{-}level insults and didn't consider the underlying emotions and power dynamics. I should have paid more attention to the fact that Luca is a chef with a strong sense of pride and ownership, while Charlie is a critic who is simply expressing his opinion. } \\
\\ 
\texttt{To improve my strategy, I will focus on identifying the key emotional triggers and power struggles in the dialogue. I will also consider the context and the characters' motivations, rather than just reacting to the insults. Specifically, I will look for moments where Luca's pride and ego are being challenged, and where Charlie's criticism is being met with defensiveness. } \\
\\ 
\texttt{I will also consider the fact that the dialogue ends with Luca storming off and Charlie standing firm, which suggests that Luca is the one who is feeling more emotionally invested and vulnerable. } \\
\\ 
\texttt{With this new understanding, I will re{-}evaluate the emotions of the characters and assign scores accordingly.} \\
\\ 
\\ 
\\ 
\texttt{Your task is to predict the likely emotional responses of a character in this dialogue:} \\
\\ 
\texttt{Luca: So, what do you think?} \\
\texttt{Charlie: It's... interesting.} \\
\texttt{Luca: Interesting? That's it?} \\
\texttt{Charlie: Well, it's not exactly a compliment, is it?} \\
\texttt{Luca: You haven't even tasted it yet. } \\
\texttt{Charlie: I don't need to taste it to know it's a mess.} \\
\texttt{Luca: A mess? You're calling my masterpiece a mess?} \\
\texttt{Charlie: If the shoe fits...} \\
\texttt{Luca: You're just a critic, Charlie. You don't know what it takes to create something.} \\
\texttt{Charlie: And you don't know what it takes to appreciate something. } \\
\texttt{Luca: I've been a chef for over 20 years. I've earned my stars.} \\
\texttt{Charlie: And I've been a critic for 15. I've earned my stripes.} \\
\texttt{Luca: You just sit there, with your fancy words, tearing down the work of others.} \\
\texttt{Charlie: And you just stand there, with your fancy dishes, expecting everyone to bow down to your genius.} \\
\texttt{Luca: You're just bitter because your life is falling apart.} \\
\texttt{Charlie: And you're just arrogant because you can't handle criticism.} \\
\texttt{Luca: I don't need to take this from you.} \\
\texttt{Charlie: And I don't need to eat this from you.} \\
\texttt{Luca: You know what, Charlie? You can keep your stripes. I don't need them.} \\
\texttt{Charlie: Good. Because I wasn't planning on giving them to you anyway.} \\
\texttt{{[}End dialogue{]}} \\
\\ 
\texttt{At the end of this dialogue, Luca would feel...} \\
\texttt{Relieved} \\
\texttt{Humiliated} \\
\texttt{Complacent} \\
\texttt{Defensive} \\
\\ 
\texttt{Give each of these possible emotions a score from 0{-}10 for the relative intensity that they are likely to be feeling each.} \\
\\ 
\texttt{You must output in the following format, including headings (of course, you should give your own scores), with no additional commentary:} \\
\\ 
\texttt{Relieved: <score>} \\
\texttt{Humiliated: <score>} \\
\texttt{Complacent: <score>} \\
\texttt{Defensive: <score>} \\
\\ 
\\ 
\texttt{{[}End of answer{]}} \\
\\ 
\texttt{Remember: zero is a valid score, meaning they are likely not feeling that emotion. You must score at least one emotion > 0.} \\
\\ 
\texttt{Your answer:} \\
            }
        }
    }
     \\ \\

    \theutterance \stepcounter{utterance}  
    & \multicolumn{4}{p{0.6\linewidth}}{
        \cellcolor[rgb]{0.561,0.737,0.561}{
            \makecell[{{p{\linewidth}}}]{
                \texttt{\tiny{[P1$\rangle$GM]}}
                \texttt{Relieved: 0} \\
\texttt{Humiliated: 9} \\
\texttt{Complacent: 0} \\
\texttt{Defensive: 8} \\
            }
        }
    }
    & & \\ \\

    \theutterance \stepcounter{utterance}  
    & & & \multicolumn{2}{p{0.3\linewidth}}{
        \cellcolor[rgb]{0.9,0.9,0.9}{
            \makecell[{{p{\linewidth}}}]{
                \texttt{\tiny{[GM$|$GM]}}
                \texttt{Relieved: 0} \\
\texttt{Humiliated: 9} \\
\texttt{Complacent: 0} \\
\texttt{Defensive: 8} \\
            }
        }
    }
    & & \\ \\

    \theutterance \stepcounter{utterance}  
    & & & \multicolumn{2}{p{0.3\linewidth}}{
        \cellcolor[rgb]{0.9,0.9,0.9}{
            \makecell[{{p{\linewidth}}}]{
                \texttt{\tiny{[GM$|$GM]}}
                \texttt{Relieved: 0} \\
\texttt{Humiliated: 6} \\
\texttt{Complacent: 0} \\
\texttt{Defensive: 8} \\
            }
        }
    }
    & & \\ \\

    \theutterance \stepcounter{utterance}  
    & & & \multicolumn{2}{p{0.3\linewidth}}{
        \cellcolor[rgb]{0.9,0.9,0.9}{
            \makecell[{{p{\linewidth}}}]{
                \texttt{\tiny{[GM$|$GM]}}
                \texttt{game\_result = LOSE} \\
            }
        }
    }
    & & \\ \\

    \theutterance \stepcounter{utterance}  
    & & \multicolumn{4}{p{0.6\linewidth}}{
        \cellcolor[rgb]{0.95,0.95,0.95}{
            \makecell[{{p{\linewidth}}}]{
                \texttt{\tiny{[P1$\leftrightarrow$P]}}
                \texttt{Post{-}episode Reflexion (not sent to GM):} \\
\texttt{New plan: } \\
\\ 
\texttt{Upon re{-}reading the dialogue, I realized that I focused too much on the surface{-}level insults and didn't consider the underlying emotions and power dynamics. I should have paid more attention to the fact that Luca is a chef with a strong sense of pride and ownership, while Charlie is a critic who is simply expressing his opinion. } \\
\\ 
\texttt{To improve my strategy, I will focus on identifying the key emotional triggers and power struggles in the dialogue. I will also consider the context and the characters' motivations, rather than just reacting to the insults. Specifically, I will look for moments where Luca's pride and ego are being challenged, and where Charlie's criticism is being met with defensiveness. } \\
\\ 
\texttt{I will also consider the fact that the dialogue ends with Luca storming off and Charlie standing firm, which suggests that Luca is the one who is feeling more emotionally invested and vulnerable. } \\
\\ 
\texttt{With this new understanding, I will re{-}evaluate the emotions of the characters and assign scores accordingly. I will pay closer attention to Luca's reactions and body language, and consider how his pride and ego are being affected by Charlie's criticism. I will also look for moments where Charlie is standing up for himself and asserting his authority as a critic. } \\
\\ 
\texttt{Based on this new understanding, I will re{-}score the emotions of the characters as follows:} \\
\\ 
\texttt{Relieved: 0} \\
\texttt{Humiliated: 10} \\
\texttt{Complacent: 0} \\
\texttt{Defensive: 9} \\
            }
        }
    }
    & \\ \\

\end{supertabular}
}

\end{document}
